%%
% Copyright (c) 2017 - 2023, Pascal Wagler;
% Copyright (c) 2014 - 2023, John MacFarlane
%
% All rights reserved.
%
% Redistribution and use in source and binary forms, with or without
% modification, are permitted provided that the following conditions
% are met:
%
% - Redistributions of source code must retain the above copyright
% notice, this list of conditions and the following disclaimer.
%
% - Redistributions in binary form must reproduce the above copyright
% notice, this list of conditions and the following disclaimer in the
% documentation and/or other materials provided with the distribution.
%
% - Neither the name of John MacFarlane nor the names of other
% contributors may be used to endorse or promote products derived
% from this software without specific prior written permission.
%
% THIS SOFTWARE IS PROVIDED BY THE COPYRIGHT HOLDERS AND CONTRIBUTORS
% "AS IS" AND ANY EXPRESS OR IMPLIED WARRANTIES, INCLUDING, BUT NOT
% LIMITED TO, THE IMPLIED WARRANTIES OF MERCHANTABILITY AND FITNESS
% FOR A PARTICULAR PURPOSE ARE DISCLAIMED. IN NO EVENT SHALL THE
% COPYRIGHT OWNER OR CONTRIBUTORS BE LIABLE FOR ANY DIRECT, INDIRECT,
% INCIDENTAL, SPECIAL, EXEMPLARY, OR CONSEQUENTIAL DAMAGES (INCLUDING,
% BUT NOT LIMITED TO, PROCUREMENT OF SUBSTITUTE GOODS OR SERVICES;
% LOSS OF USE, DATA, OR PROFITS; OR BUSINESS INTERRUPTION) HOWEVER
% CAUSED AND ON ANY THEORY OF LIABILITY, WHETHER IN CONTRACT, STRICT
% LIABILITY, OR TORT (INCLUDING NEGLIGENCE OR OTHERWISE) ARISING IN
% ANY WAY OUT OF THE USE OF THIS SOFTWARE, EVEN IF ADVISED OF THE
% POSSIBILITY OF SUCH DAMAGE.
%%

%%
% This is the Eisvogel pandoc LaTeX template.
%
% For usage information and examples visit the official GitHub page:
% https://github.com/Wandmalfarbe/pandoc-latex-template
%%

% Options for packages loaded elsewhere
\PassOptionsToPackage{unicode}{hyperref}
\PassOptionsToPackage{hyphens}{url}
\PassOptionsToPackage{dvipsnames,svgnames,x11names,table}{xcolor}
%
\documentclass[
  paper=a4,
  ,captions=tableheading
]{scrartcl}
\usepackage{amsmath,amssymb}
% Use setspace anyway because we change the default line spacing.
% The spacing is changed early to affect the titlepage and the TOC.
\usepackage{setspace}
\setstretch{1.2}
\usepackage{iftex}
\ifPDFTeX
  \usepackage[T1]{fontenc}
  \usepackage[utf8]{inputenc}
  \usepackage{textcomp} % provide euro and other symbols
\else % if luatex or xetex
  \usepackage{unicode-math} % this also loads fontspec
  \defaultfontfeatures{Scale=MatchLowercase}
  \defaultfontfeatures[\rmfamily]{Ligatures=TeX,Scale=1}
\fi
\usepackage{lmodern}
\ifPDFTeX\else
  % xetex/luatex font selection
\fi
% Use upquote if available, for straight quotes in verbatim environments
\IfFileExists{upquote.sty}{\usepackage{upquote}}{}
\IfFileExists{microtype.sty}{% use microtype if available
  \usepackage[]{microtype}
  \UseMicrotypeSet[protrusion]{basicmath} % disable protrusion for tt fonts
}{}
\makeatletter
\@ifundefined{KOMAClassName}{% if non-KOMA class
  \IfFileExists{parskip.sty}{%
    \usepackage{parskip}
  }{% else
    \setlength{\parindent}{0pt}
    \setlength{\parskip}{6pt plus 2pt minus 1pt}}
}{% if KOMA class
  \KOMAoptions{parskip=half}}
\makeatother
\usepackage{xcolor}
\definecolor{default-linkcolor}{HTML}{A50000}
\definecolor{default-filecolor}{HTML}{A50000}
\definecolor{default-citecolor}{HTML}{4077C0}
\definecolor{default-urlcolor}{HTML}{4077C0}
\usepackage[margin=2.5cm,includehead=true,includefoot=true,centering,]{geometry}
\usepackage{color}
\usepackage{fancyvrb}
\newcommand{\VerbBar}{|}
\newcommand{\VERB}{\Verb[commandchars=\\\{\}]}
\DefineVerbatimEnvironment{Highlighting}{Verbatim}{commandchars=\\\{\}}
% Add ',fontsize=\small' for more characters per line
\newenvironment{Shaded}{}{}
\newcommand{\AlertTok}[1]{\textcolor[rgb]{1.00,0.00,0.00}{\textbf{#1}}}
\newcommand{\AnnotationTok}[1]{\textcolor[rgb]{0.38,0.63,0.69}{\textbf{\textit{#1}}}}
\newcommand{\AttributeTok}[1]{\textcolor[rgb]{0.49,0.56,0.16}{#1}}
\newcommand{\BaseNTok}[1]{\textcolor[rgb]{0.25,0.63,0.44}{#1}}
\newcommand{\BuiltInTok}[1]{\textcolor[rgb]{0.00,0.50,0.00}{#1}}
\newcommand{\CharTok}[1]{\textcolor[rgb]{0.25,0.44,0.63}{#1}}
\newcommand{\CommentTok}[1]{\textcolor[rgb]{0.38,0.63,0.69}{\textit{#1}}}
\newcommand{\CommentVarTok}[1]{\textcolor[rgb]{0.38,0.63,0.69}{\textbf{\textit{#1}}}}
\newcommand{\ConstantTok}[1]{\textcolor[rgb]{0.53,0.00,0.00}{#1}}
\newcommand{\ControlFlowTok}[1]{\textcolor[rgb]{0.00,0.44,0.13}{\textbf{#1}}}
\newcommand{\DataTypeTok}[1]{\textcolor[rgb]{0.56,0.13,0.00}{#1}}
\newcommand{\DecValTok}[1]{\textcolor[rgb]{0.25,0.63,0.44}{#1}}
\newcommand{\DocumentationTok}[1]{\textcolor[rgb]{0.73,0.13,0.13}{\textit{#1}}}
\newcommand{\ErrorTok}[1]{\textcolor[rgb]{1.00,0.00,0.00}{\textbf{#1}}}
\newcommand{\ExtensionTok}[1]{#1}
\newcommand{\FloatTok}[1]{\textcolor[rgb]{0.25,0.63,0.44}{#1}}
\newcommand{\FunctionTok}[1]{\textcolor[rgb]{0.02,0.16,0.49}{#1}}
\newcommand{\ImportTok}[1]{\textcolor[rgb]{0.00,0.50,0.00}{\textbf{#1}}}
\newcommand{\InformationTok}[1]{\textcolor[rgb]{0.38,0.63,0.69}{\textbf{\textit{#1}}}}
\newcommand{\KeywordTok}[1]{\textcolor[rgb]{0.00,0.44,0.13}{\textbf{#1}}}
\newcommand{\NormalTok}[1]{#1}
\newcommand{\OperatorTok}[1]{\textcolor[rgb]{0.40,0.40,0.40}{#1}}
\newcommand{\OtherTok}[1]{\textcolor[rgb]{0.00,0.44,0.13}{#1}}
\newcommand{\PreprocessorTok}[1]{\textcolor[rgb]{0.74,0.48,0.00}{#1}}
\newcommand{\RegionMarkerTok}[1]{#1}
\newcommand{\SpecialCharTok}[1]{\textcolor[rgb]{0.25,0.44,0.63}{#1}}
\newcommand{\SpecialStringTok}[1]{\textcolor[rgb]{0.73,0.40,0.53}{#1}}
\newcommand{\StringTok}[1]{\textcolor[rgb]{0.25,0.44,0.63}{#1}}
\newcommand{\VariableTok}[1]{\textcolor[rgb]{0.10,0.09,0.49}{#1}}
\newcommand{\VerbatimStringTok}[1]{\textcolor[rgb]{0.25,0.44,0.63}{#1}}
\newcommand{\WarningTok}[1]{\textcolor[rgb]{0.38,0.63,0.69}{\textbf{\textit{#1}}}}

% Workaround/bugfix from jannick0.
% See https://github.com/jgm/pandoc/issues/4302#issuecomment-360669013)
% or https://github.com/Wandmalfarbe/pandoc-latex-template/issues/2
%
% Redefine the verbatim environment 'Highlighting' to break long lines (with
% the help of fvextra). Redefinition is necessary because it is unlikely that
% pandoc includes fvextra in the default template.
\usepackage{fvextra}
\DefineVerbatimEnvironment{Highlighting}{Verbatim}{breaklines,fontsize=\small,commandchars=\\\{\}}

\usepackage{longtable,booktabs,array}
\usepackage{calc} % for calculating minipage widths
% Correct order of tables after \paragraph or \subparagraph
\usepackage{etoolbox}
\makeatletter
\patchcmd\longtable{\par}{\if@noskipsec\mbox{}\fi\par}{}{}
\makeatother
% Allow footnotes in longtable head/foot
\IfFileExists{footnotehyper.sty}{\usepackage{footnotehyper}}{\usepackage{footnote}}
\makesavenoteenv{longtable}
% add backlinks to footnote references, cf. https://tex.stackexchange.com/questions/302266/make-footnote-clickable-both-ways
\usepackage{footnotebackref}
\setlength{\emergencystretch}{3em} % prevent overfull lines
\providecommand{\tightlist}{%
  \setlength{\itemsep}{0pt}\setlength{\parskip}{0pt}}
\setcounter{secnumdepth}{-\maxdimen} % remove section numbering
\ifLuaTeX
  \usepackage{selnolig}  % disable illegal ligatures
\fi
\IfFileExists{bookmark.sty}{\usepackage{bookmark}}{\usepackage{hyperref}}
\IfFileExists{xurl.sty}{\usepackage{xurl}}{} % add URL line breaks if available
\urlstyle{same}
\hypersetup{
  pdftitle={PasswordStore Vulnerability Assessment},
  pdfauthor={Cyfrin.io-Student-Elechukwu},
  hidelinks,
  breaklinks=true,
  pdfcreator={LaTeX via pandoc with the Eisvogel template}}
\title{PasswordStore Vulnerability Assessment}
\usepackage{etoolbox}
\makeatletter
\providecommand{\subtitle}[1]{% add subtitle to \maketitle
  \apptocmd{\@title}{\par {\large #1 \par}}{}{}
}
\makeatother
\subtitle{Protocol Audit Report}
\author{Cyfrin.io-Student-Elechukwu}
\date{January 22, 2026}



%%
%% added
%%


%
% for the background color of the title page
%

%
% break urls
%
\PassOptionsToPackage{hyphens}{url}

%
% When using babel or polyglossia with biblatex, loading csquotes is recommended
% to ensure that quoted texts are typeset according to the rules of your main language.
%
\usepackage{csquotes}

%
% captions
%
\definecolor{caption-color}{HTML}{777777}
\usepackage[font={stretch=1.2}, textfont={color=caption-color}, position=top, skip=4mm, labelfont=bf, singlelinecheck=false, justification=raggedright]{caption}
\setcapindent{0em}

%
% blockquote
%
\definecolor{blockquote-border}{RGB}{221,221,221}
\definecolor{blockquote-text}{RGB}{119,119,119}
\usepackage{mdframed}
\newmdenv[rightline=false,bottomline=false,topline=false,linewidth=3pt,linecolor=blockquote-border,skipabove=\parskip]{customblockquote}
\renewenvironment{quote}{\begin{customblockquote}\list{}{\rightmargin=0em\leftmargin=0em}%
\item\relax\color{blockquote-text}\ignorespaces}{\unskip\unskip\endlist\end{customblockquote}}

%
% Source Sans Pro as the de­fault font fam­ily
% Source Code Pro for monospace text
%
% 'default' option sets the default
% font family to Source Sans Pro, not \sfdefault.
%
\ifnum 0\ifxetex 1\fi\ifluatex 1\fi=0 % if pdftex
    \usepackage[default]{sourcesanspro}
  \usepackage{sourcecodepro}
  \else % if not pdftex
    \usepackage[default]{sourcesanspro}
  \usepackage{sourcecodepro}

  % XeLaTeX specific adjustments for straight quotes: https://tex.stackexchange.com/a/354887
  % This issue is already fixed (see https://github.com/silkeh/latex-sourcecodepro/pull/5) but the
  % fix is still unreleased.
  % TODO: Remove this workaround when the new version of sourcecodepro is released on CTAN.
  \ifxetex
    \makeatletter
    \defaultfontfeatures[\ttfamily]
      { Numbers   = \sourcecodepro@figurestyle,
        Scale     = \SourceCodePro@scale,
        Extension = .otf }
    \setmonofont
      [ UprightFont    = *-\sourcecodepro@regstyle,
        ItalicFont     = *-\sourcecodepro@regstyle It,
        BoldFont       = *-\sourcecodepro@boldstyle,
        BoldItalicFont = *-\sourcecodepro@boldstyle It ]
      {SourceCodePro}
    \makeatother
  \fi
  \fi

%
% heading color
%
\definecolor{heading-color}{RGB}{40,40,40}
\addtokomafont{section}{\color{heading-color}}
% When using the classes report, scrreprt, book,
% scrbook or memoir, uncomment the following line.
%\addtokomafont{chapter}{\color{heading-color}}

%
% variables for title, author and date
%
\usepackage{titling}
\title{PasswordStore Vulnerability Assessment}
\author{Cyfrin.io-Student-Elechukwu}
\date{January 22, 2026}

%
% tables
%

\definecolor{table-row-color}{HTML}{F5F5F5}
\definecolor{table-rule-color}{HTML}{999999}

%\arrayrulecolor{black!40}
\arrayrulecolor{table-rule-color}     % color of \toprule, \midrule, \bottomrule
\setlength\heavyrulewidth{0.3ex}      % thickness of \toprule, \bottomrule
\renewcommand{\arraystretch}{1.3}     % spacing (padding)


%
% remove paragraph indention
%
\setlength{\parindent}{0pt}
\setlength{\parskip}{6pt plus 2pt minus 1pt}
\setlength{\emergencystretch}{3em}  % prevent overfull lines

%
%
% Listings
%
%


%
% header and footer
%
\usepackage[headsepline,footsepline]{scrlayer-scrpage}

\newpairofpagestyles{eisvogel-header-footer}{
  \clearpairofpagestyles
  \ihead*{PasswordStore Vulnerability Assessment}
  \chead*{}
  \ohead*{January 22, 2026}
  \ifoot*{Cyfrin.io-Student-Elechukwu}
  \cfoot*{}
  \ofoot*{\thepage}
  \addtokomafont{pageheadfoot}{\upshape}
}
\pagestyle{eisvogel-header-footer}



%%
%% end added
%%

\begin{document}

%%
%% begin titlepage
%%

%%
%% end titlepage
%%

% \maketitle


Prepared by:
{[}Cyfrin{]}{[}Elechukwu{]}(https://cyfrin.io)(https://github.com/Azel-C/Blockchain-Training)
Lead Auditors: - xxxxxxx

\section{Table of Contents}\label{table-of-contents}

\begin{itemize}
\tightlist
\item
  \hyperref[table-of-contents]{Table of Contents}
\item
  \hyperref[protocol-summary]{Protocol Summary}
\item
  \hyperref[disclaimer]{Disclaimer}
\item
  \hyperref[risk-classification]{Risk Classification}
\item
  \hyperref[audit-details]{Audit Details}

  \begin{itemize}
  \tightlist
  \item
    \hyperref[scope]{Scope}
  \item
    \hyperref[roles]{Roles}
  \end{itemize}
\item
  \hyperref[executive-summary]{Executive Summary}

  \begin{itemize}
  \tightlist
  \item
    \hyperref[issues-found]{Issues found}
  \end{itemize}
\item
  \hyperref[findings]{Findings}

  \begin{itemize}
  \tightlist
  \item
    \hyperref[high]{High}

    \begin{itemize}
    \tightlist
    \item
      \hyperref[h-1-storing-the-password-on-chain-makes-it-visible-to-anyone-and-no-longer-private]{{[}H-1{]}
      Storing the password on-chain makes it visible to anyone, and no
      longer private}
    \item
      \hyperref[h-2-passwordstoresetpassword-has-no-access-controls-meaning-non-owner-could-change-the-password]{{[}H-2{]}
      \texttt{PasswordStore::setPassword} has no access controls,
      meaning non-owner could change the password.}
    \end{itemize}
  \item
    \hyperref[informational]{Informational}

    \begin{itemize}
    \tightlist
    \item
      \hyperref[i-1-the-passwordstoresetpassword-natspec-indicates-a-parameter-that-doesnt-exist-causing-the-natspec-to-be-incorrect]{{[}I-1{]}
      The \texttt{PasswordStore::setPassword} natspec indicates a
      parameter that doesn't exist, causing the natspec to be
      incorrect.}
    \end{itemize}
  \end{itemize}
\end{itemize}

\section{Protocol Summary}\label{protocol-summary}

PasswordStore is a protocol dedicated to storage and retrival of user's
passwords. The protocol is designed to be used by a single user, and not
designed to be used by multiple user's. Only the owner should be able to
set and access the password.

\section{Disclaimer}\label{disclaimer}

The Elechukwu team makes all effort to find as many vulnerabilities in
the code in the given time period, but holds no responsibilities for the
findings provided in this document. A security audit by the team is not
an endorsement of the underlying business or product. The audit was
time-boxed and the review of the code was solely on the security aspects
of the Solidity implementation of the contracts.

\section{Risk Classification}\label{risk-classification}

{\def\LTcaptype{none} % do not increment counter
\begin{longtable}[]{@{}lllll@{}}
\toprule\noalign{}
& & Impact & & \\
\midrule\noalign{}
\endhead
\bottomrule\noalign{}
\endlastfoot
& & High & Medium & Low \\
& High & H & H/M & M \\
Likelihood & Medium & H/M & M & M/L \\
& Low & M & M/L & L \\
\end{longtable}
}

We use the
\href{https://docs.codehawks.com/hawks-auditors/how-to-evaluate-a-finding-severity}{CodeHawks}
severity matrix to determine severity. See the documentation for more
details.

\section{Audit Details}\label{audit-details}

\textbf{The findings described in this document correspond the following
commit hash:}

\begin{verbatim}
    2e8f81e263b3a9d18fab4fb5c46805ffc10a9990
\end{verbatim}

\subsection{Scope}\label{scope}

\begin{verbatim}
./src/
#-- PasswordStore.sol
\end{verbatim}

\subsection{Roles}\label{roles}

\begin{itemize}
\tightlist
\item
  Owner: The user who can set the password and read the password.
\item
  Outsider: No one else should be able to read or update the password.
\end{itemize}

\section{Executive Summary}\label{executive-summary}

\subsection{Issues found}\label{issues-found}

{\def\LTcaptype{none} % do not increment counter
\begin{longtable}[]{@{}ll@{}}
\toprule\noalign{}
Severity & Number of issues found \\
\midrule\noalign{}
\endhead
\bottomrule\noalign{}
\endlastfoot
Highs & 2 \\
Mediums & 0 \\
Lows & 0 \\
Info & 1 \\
Total & 3 \\
\end{longtable}
}

\section{Findings}\label{findings}

\subsection{High}\label{high}

\subsubsection{{[}H-1{]} Storing the password on-chain makes it visible
to anyone, and no longer
private}\label{h-1-storing-the-password-on-chain-makes-it-visible-to-anyone-and-no-longer-private}

\textbf{Description:} All data stored on-chain is visible to anyone, and
can be read directly from the blockchain. The
\texttt{PasswordStore::s\_password} varible is intended to be a private
varible and only accessed through the
\texttt{PasswordStore::getPassword} function, which is intened to be
only called by the owner of the contract.

We show one such method of reading data off chain below.

\textbf{Impact:} Anyone can read the private password, which severly
break the functionality of the protocol.

\textbf{Proof of Concept:} (Proof of Code)

The below test case shows anyone can read the password directly from the
blockchain.

\begin{enumerate}
\def\labelenumi{\arabic{enumi}.}
\tightlist
\item
  Create a locally runing chain
\end{enumerate}

\begin{Shaded}
\begin{Highlighting}[]
\FunctionTok{make}\NormalTok{ Anvil}
\end{Highlighting}
\end{Shaded}

\begin{enumerate}
\def\labelenumi{\arabic{enumi}.}
\setcounter{enumi}{1}
\tightlist
\item
  Deploy the contract to chain
\end{enumerate}

\begin{verbatim}
make deploy
\end{verbatim}

\begin{enumerate}
\def\labelenumi{\arabic{enumi}.}
\setcounter{enumi}{2}
\tightlist
\item
  Run the storage tool we use \texttt{1} because that's the stroage slot
  of \texttt{s\_password} in the contract.
\end{enumerate}

\begin{verbatim}
cast storage <ADDRESS_HERE> 1 --rpc-url http://127.0.0.1:8545
\end{verbatim}

you'll get an output like this

\texttt{0x6d7950617373776f726400000000000000000000000000000000000000000014}

you can then parse the hex to a string, like this:

\begin{verbatim}
cast parse-bytes32-string 0x6d7950617373776f726400000000000000000000000000000000000000000014
\end{verbatim}

And get an output of:

\begin{verbatim}
myPassword
\end{verbatim}

\textbf{Recommended Mitigation:} Due to this, the overall architecture
of the contract should be rethought. one could encrypt the password
of-chain, and then strore the encrypted password on-chain. this will
require the user to remember another password off-chain to decrypt the
password. However, you'd also want to remove the view function, as you
wouldn't want the user to accidentally send transaction with the
password that decrypts your password.

\subsubsection{\texorpdfstring{{[}H-2{]}
\texttt{PasswordStore::setPassword} has no access controls, meaning
non-owner could change the
password.}{{[}H-2{]} PasswordStore::setPassword has no access controls, meaning non-owner could change the password.}}\label{h-2-passwordstoresetpassword-has-no-access-controls-meaning-non-owner-could-change-the-password.}

\textbf{Description:} The \texttt{PasswordStore::setPassword} function
is set to be an external function, howerver, the natspec of the function
and overall purpose of the Smart Contract is that
\texttt{This\ function\ allows\ only\ the\ owner\ to\ set\ a\ new\ password.}

\begin{Shaded}
\begin{Highlighting}[]
\KeywordTok{function} \FunctionTok{setPassword}\NormalTok{(string memory newPassword) external \{}
\NormalTok{@}\OperatorTok{\textgreater{}}  \CommentTok{// @audit: There is no access controls}
\NormalTok{        s\_password }\OperatorTok{=}\NormalTok{ newPassword}\OperatorTok{;}
\NormalTok{        emit }\FunctionTok{SetNetPassword}\NormalTok{()}\OperatorTok{;}
\NormalTok{    \}}
\end{Highlighting}
\end{Shaded}

\textbf{Impact:} Anyone can set/change the password of the contract,
severly breaking the contract intended functionality.

\textbf{Proof of Concept:} Add the following to
\texttt{PasswordStore.t.sol} test file.

Code

\begin{Shaded}
\begin{Highlighting}[]
\KeywordTok{function} \FunctionTok{test\_anyone\_can\_set\_password}\NormalTok{(address randomAddress) }\KeywordTok{public}\NormalTok{ \{}
\NormalTok{        vm}\OperatorTok{.}\FunctionTok{assume}\NormalTok{(randomAddress }\OperatorTok{!=}\NormalTok{ owner)}\OperatorTok{;}
\NormalTok{        vm}\OperatorTok{.}\FunctionTok{prank}\NormalTok{(randomAddress)}\OperatorTok{;}
\NormalTok{        string memory newPassword }\OperatorTok{=} \StringTok{"hackedPassword"}\OperatorTok{;}
\NormalTok{        passwordStore}\OperatorTok{.}\FunctionTok{setPassword}\NormalTok{(newPassword)}\OperatorTok{;}

\NormalTok{        vm}\OperatorTok{.}\FunctionTok{prank}\NormalTok{(owner)}\OperatorTok{;}
\NormalTok{        string memory actualPassword }\OperatorTok{=}\NormalTok{ passwordStore}\OperatorTok{.}\FunctionTok{getPassword}\NormalTok{()}\OperatorTok{;}
        \FunctionTok{assertEq}\NormalTok{(actualPassword}\OperatorTok{,}\NormalTok{ newPassword)}\OperatorTok{;}
\NormalTok{    \}}
\end{Highlighting}
\end{Shaded}

\textbf{Recommended Mitigation:} Add an access control conditional to
the \texttt{setPassword} function.

\begin{Shaded}
\begin{Highlighting}[]

\ControlFlowTok{if}\NormalTok{(msg}\OperatorTok{.}\AttributeTok{sender} \OperatorTok{!=}\NormalTok{ s\_owner)\{}
\NormalTok{  revert PasswordStore\_NotOwner}
\NormalTok{\}}
\end{Highlighting}
\end{Shaded}

\subsection{Informational}\label{informational}

\#\#\# {[}I-1{]} The \texttt{PasswordStore::setPassword} natspec
indicates a parameter that doesn't exist, causing the natspec to be
incorrect.

\textbf{Description:}

\begin{Shaded}
\begin{Highlighting}[]
    \CommentTok{/*}
\CommentTok{     * @notice This allows only the owner to retrieve the password.}
\CommentTok{     * @param newPassword The new password to set.}
\CommentTok{     */}
    \KeywordTok{function} \FunctionTok{getPassword}\NormalTok{() external view }\FunctionTok{returns}\NormalTok{ (string memory) \{}
\end{Highlighting}
\end{Shaded}

The \texttt{PasswordStore::getPassword} function signature is
\texttt{getPassword} which the natspec say it should be
\texttt{getPassword(string)}.

\textbf{Impact:} The natspec is incorrect.

\textbf{Recommended Mitigation:} Remove the incorrect natspec line.

\begin{Shaded}
\begin{Highlighting}[]
\StringTok{{-}  * @param newPassword The new password to set.}
\end{Highlighting}
\end{Shaded}


\end{document}